% Generated by scripts/build-paper.py — edit paper/src/survey.src.md
\section{Scope, definitions and what is documented}\label{sec:scope-definitions-and-what-is-documented}

\subsection{The interface}

A caller sends a \texttt{state} (a string, JSON object or array of text values) and a map of typed questions; the model evaluates every question independently against the shared state and returns one answer per question \citep{typesafe2026state,typesafe2026primitives}. Three primitives exist. A \textbf{Choice} selects among named options, each defined by a description; the answer includes the chosen option, a probability for every option and a \texttt{confidence} value. The launch post states support for up to 255 options, with a two-stage procedure for high cardinality \citep{typesafe2026launch}. A \textbf{Score} places the answer on two to ten ordered, described levels; it returns a position that may fall between levels, probabilities per level and \texttt{confidence}. The vendor warns against interpolating exact magnitudes from it \citep{typesafe2026jaggedness}. A \textbf{Noul} returns the probability that a yes/no proposition holds and carries no separate confidence \citep{typesafe2026primitives,typesafe2026confidence}.

At the cutoff the documented model was \texttt{jev-1.13.0}. The aliases \texttt{jev-latest} and \texttt{jev-preview} both pointed to it, and the vendor advises pinning the versioned ID when thresholds are tuned to a version \citep{typesafe2026models}. Input is text only; images, audio and video must first be converted to text or structured fields, and English is the primary training language \citep{typesafe2026models,typesafe2026state}. Visual Jev and JEVQA, discussed below, do not change this: one is an independent model and the other feeds text features to hosted Jev.

\subsection{Probability semantics}

Probabilities and confidence are different objects. The documentation defines \texttt{confidence} as a statistic computed from the returned distribution and recommends thresholds tested on the user's own data \citep{typesafe2026confidence}. The exact statistic is not specified in prose; an interactive demo on the documentation page computes a normalised maximum probability and labels it as how \emph{that demo} computes confidence. We therefore do not treat the demo formula as the production implementation. Empirically, one study found that native confidence ranked items almost exactly like the maximum label probability (Spearman 0.948–0.999) \citep{arxiv2609_26550}. The vendor also states that calibration is measured across groups of predictions and does not guarantee that an individual answer is correct \citep{typesafe2026systemone}. Two independent sources observe that probabilities are returned on a 0.01 grid, which bounds calibration error for rare events from below and can put exactly zero on a correct answer \citep{arxiv2609_24052,gh_scienthoon_jev_ood_calibration}.

\subsection{What is not disclosed}

The launch post names a new architecture, a parallel sampler and the RLCD training method, and the documentation says one set of weights serves every account without per-customer fine-tuning \citep{typesafe2026launch,typesafe2026models}. The architecture, training data and objective are not published. Community reconstructions exist; Kev, for example, follows an architecture described in a third-party blog post. They remain speculation and are not evidence about the commercial model \citep{gh_jaredpalmer_kev}.

\subsection{Vendor claims as hypotheses}

The headline “193.6× faster, 444.6× cheaper” comes from the vendor's workflow evaluations. Their reference answers average two external frontier models; the vendor expects the gains to be at the high end of real-world use and notes that its own capabilities team wrote the workflows \citep{typesafe2026launch}. The plotted 0\% type-error rate is described as analytical rather than empirical, and latency figures were measured from the vendor's West Coast laptops \citep{typesafe2026launch}. We record these as vendor claims, not as evaluations.

\subsection{Names that collide}

Several names mean different things. TypeSafe's \emph{System One} borrows Kahneman's System 1 metaphor, but System 1/2 reasoning literature predates the product and is broader \citep{arxiv2502_17419}. \emph{RLCD} in the Jev context expands to Reinforcement Learning for Calibrated Decisions and is unrelated to RLCD, Reinforcement Learning from Contrastive Distillation (2023); one core study makes the same distinction \citep{arxiv2609_24574}. \emph{OpenJev} names three different things: the former name of SemIf (a frozen-decoder readout), a DiffusionGemma server, and the “Open-Jev” of one paper's title (JevLite) \citep{gh_theoleecj_semif,gh_razorback16_openjev,arxiv2609_23959}. \emph{Visual Jev} is independent research, not a TypeSafe release \citep{arxiv2609_25845}. One repository called a “synthetic survey” is a questionnaire experiment \citep{gh_jjd_lab_jev_synthetic_survey}. Adjacency in name implies neither the same algorithm nor copying.
